\subsection{Related work and research positioning}
\label{sec:related_work}

The previous chapter introduced the core concepts needed for the project.
This section narrows the focus to the literature most directly relevant to the present study and clarifies what the project contributes relative to that literature.
The aim is not to survey all of adversarial machine learning or all of explainable AI, but to identify the specific strands of work that motivate the chosen experimental design.

\subsection{Sparse and one-pixel adversarial attacks}
A large proportion of adversarial machine learning literature studies perturbations under norm-bounded threat models, especially $\ell_\infty$ and $\ell_2$ settings \cite{goodfellow2015explaining,madry2018towards}.
That literature is important because it establishes the general fact that modern neural networks can be systematically manipulated by carefully chosen inputs.
However, the one-pixel setting is more specific.
Rather than asking how much bounded noise can be added overall, it asks whether a classifier can be changed by modifying only one spatial location.
This makes the perturbation extremely sparse and turns the attack problem into a compact search over a small number of continuous variables.

The one-pixel attack of Su et al.\ is therefore notable not simply because it is visually striking, but because it provides an unusually interpretable black-box optimisation problem \cite{su2019onepixel}.
The attacker is forced to concentrate all adversarial leverage into a single coordinate and three colour values.
That is scientifically useful because success under such a restricted budget reveals how locally fragile a learned decision boundary can be.
It is also methodologically useful because the restricted search space makes the trade-off between attack success, search strategy, and query cost easier to study than in higher-dimensional adversarial settings.

The broader adversarial literature also includes work on perturbations that generalise across inputs, such as universal adversarial perturbations \cite{moosavi2017universal}.
Although such work is not directly about one-pixel attacks, it reinforces the point that adversarial behaviour should not be treated as an isolated curiosity tied to one optimisation trick.
The present report operates at the opposite end of the perturbation spectrum: instead of a broad perturbation pattern that transfers widely, it studies an extreme sparse-budget case in which the perturbation is highly local and individually optimised.

\subsection{Black-box attack strategies and query efficiency}
Once the white-box assumption is removed, adversarial search becomes substantially harder.
The black-box literature has therefore explored several broad strategies, including transfer-based attacks, query-based zeroth-order optimisation, and decision-based search.
Papernot et al.\ showed that black-box attacks can be mounted by training substitute models and exploiting transferability \cite{papernot2017practical}.
Chen et al.\ introduced ZOO, which approximates gradient information through zeroth-order optimisation without requiring a substitute model \cite{chen2017zoo}.
Brendel et al.\ studied decision-based attacks that require only model decisions rather than rich score information \cite{brendel2018decision}.
Guo et al.\ further emphasised the practical importance of simple query-based black-box attacks \cite{guo2019simplebb}.

What these works collectively show is that success rate alone is an incomplete way to evaluate a black-box attack.
Query cost matters because queries are the scarce resource in realistic black-box settings.
An attack that succeeds eventually but requires substantially more interaction may be much less attractive in practice than one that achieves a similar success rate with fewer evaluations.
This point is especially important for the present project because RQ2 is explicitly about the success--cost trade-off, not merely about whether guidance can preserve attack success.

The present report differs from this broader black-box literature in two ways.
First, it operates in the extremely sparse one-pixel regime rather than in a more flexible perturbation space.
Second, it evaluates guidance using a black-box saliency prior rather than transferability or surrogate gradients.
The project therefore does not compete with the strongest general-purpose black-box attacks in the literature.
Instead, it studies whether explanation-derived spatial guidance is useful within a tightly controlled sparse-budget setting.

\subsection{Guidance priors for adversarial search}
Much of the literature on stronger adversarial attacks improves optimisation by exploiting gradient information, surrogate models, or carefully designed objectives \cite{madry2018towards,chen2017zoo,guo2019simplebb}.
By contrast, the present project asks whether a post-hoc explanation can act as a practical prior over where a sparse black-box search should focus.
This is a more specific question than ``how can black-box attacks be made stronger in general?''
It is instead the question of whether an explanation-derived estimate of current importance is a helpful spatial bias for one-pixel search.

This distinction matters because a saliency prior is not equivalent to gradient access.
A saliency map explains, approximately and post hoc, which parts of the current input appear influential for the model's current prediction.
That is not the same as the most query-efficient direction for causing misclassification.
A guidance prior may still be useful, but its value is empirical rather than guaranteed.
The literature therefore motivates the idea, but does not settle it.

That gap is part of the reason the present project is worth doing.
The project does not claim to introduce a novel adversarial optimiser at the level of the broader black-box literature.
Its narrower contribution is to test whether explanation-derived guidance has measurable value in a sparse black-box one-pixel regime, and to evaluate that value in terms of both attack success and query cost.

\subsection{Adversarial perturbations and post-hoc explanations}
The second major literature strand concerns post-hoc explanation methods and their robustness.
Grad-CAM, Integrated Gradients, and RISE represent three important explanation paradigms: activation-based localisation, gradient-based attribution, and black-box perturbation-based attribution respectively \cite{selvaraju2017gradcam,sundararajan2017ig,petsiuk2018rise}.
These methods are widely used, but the literature has repeatedly shown that explanation outputs should not be treated as automatically reliable or stable.

Adebayo et al.\ demonstrated that some saliency methods can fail important sanity checks, showing that visually plausible explanation maps may not be as informative as they appear \cite{adebayo2018sanity}.
Work on explanation fragility and manipulation has pushed this further.
Ghorbani et al.\ showed that neural network interpretations can be fragile under adversarial perturbation \cite{ghorbani2019fragile}.
Alvarez-Melis and Jaakkola studied the robustness of interpretability methods more directly, emphasising that interpretability outputs themselves may vary sharply under small input changes \cite{alvarezmelis2018robustness}.
Dombrowski et al.\ showed that explanations can be manipulated in ways driven by geometric properties rather than faithful semantic reasoning \cite{dombrowski2019explanations}.

This literature is highly relevant to RQ1.
If explanation methods can be fragile or manipulable, then it should not be assumed that a successful one-pixel adversarial example leaves explanation structure intact.
At the same time, much of the explanation-robustness literature studies explanation instability as a phenomenon in its own right.
The present project instead embeds explanation analysis within a one-pixel attack pipeline and compares multiple attack variants under the same evaluation regime.
That makes explanation robustness part of a larger comparative study rather than a standalone explanation paper.

\subsection{Research positioning and contribution boundary}
The project should be positioned as a reproducible experimental study that combines three strands of literature: sparse black-box attacks, post-hoc explanation robustness, and guidance-based attack design.
Its contribution is not a new explanation method, not a new state-of-the-art black-box attack, and not a universal robustness claim.
The contribution is instead the combination itself: a controlled 10k-scale evaluation in which one-pixel attack success, query cost, explanation stability, and deletion-based faithfulness are analysed together.

Relative to Su et al., the project revisits the one-pixel setting with a stronger emphasis on reproducible large-scale comparison and explanation analysis \cite{su2019onepixel}.
Relative to the broader black-box attack literature, the project studies a narrower perturbation setting but asks a different question: whether explanation-derived guidance helps in a sparse regime.
Relative to the explanation-robustness literature, the project does not attempt to characterise explanation fragility in full generality, but instead asks how explanation behaviour changes under successful one-pixel attacks and how that behaviour varies across attack procedures.

The contribution boundary is therefore deliberately modest but clear.
The report claims evidence about a fixed dataset, a fixed pretrained architecture, three explanation methods, and three attack configurations under a single-pixel black-box threat model.
Within that scope, the project contributes a more integrated and auditable account of the relationship between sparse adversarial success, attack cost, and explanation behaviour than is typically found in a compact one-pixel reproduction study.
